\documentclass[10pt,a4paper]{amsart}
\usepackage{amsmath,amssymb,amsthm,graphicx}
% Retained minimal preamble for the Stacks Project chapter more-morphisms.tex (More on Morphisms).
% Verbatim excerpt of the chapter preamble: Xy-pic setup (native preamble lines 14-18) and the native
% theorem declarations with the mathematical macros (native preamble lines 154-217). Copied mechanically;
% nothing added, removed or redefined. The wrapper class and the amsmath/amssymb/amsthm/graphicx base
% packages are supplied by the assembler, so the class/externaldocument part of the native preamble is omitted.
% For commutative diagrams we use Xy-pic
\usepackage[all]{xy}

% We use 2cell for 2-commutative diagrams.
\xyoption{2cell}
% Theorem environments.
%
\theoremstyle{plain}
\newtheorem{theorem}[subsection]{Theorem}
\newtheorem{proposition}[subsection]{Proposition}
\newtheorem{lemma}[subsection]{Lemma}

\theoremstyle{definition}
\newtheorem{definition}[subsection]{Definition}
\newtheorem{example}[subsection]{Example}
\newtheorem{exercise}[subsection]{Exercise}
\newtheorem{situation}[subsection]{Situation}

\theoremstyle{remark}
\newtheorem{remark}[subsection]{Remark}
\newtheorem{remarks}[subsection]{Remarks}

\numberwithin{equation}{subsection}

% Macros
%
\def\lim{\mathop{\mathrm{lim}}\nolimits}
\def\colim{\mathop{\mathrm{colim}}\nolimits}
\def\Spec{\mathop{\mathrm{Spec}}}
\def\Hom{\mathop{\mathrm{Hom}}\nolimits}
\def\Ext{\mathop{\mathrm{Ext}}\nolimits}
\def\SheafHom{\mathop{\mathcal{H}\!\mathit{om}}\nolimits}
\def\SheafExt{\mathop{\mathcal{E}\!\mathit{xt}}\nolimits}
\def\Sch{\mathit{Sch}}
\def\Mor{\mathop{\mathrm{Mor}}\nolimits}
\def\Ob{\mathop{\mathrm{Ob}}\nolimits}
\def\Sh{\mathop{\mathit{Sh}}\nolimits}
\def\NL{\mathop{N\!L}\nolimits}
\def\CH{\mathop{\mathrm{CH}}\nolimits}
\def\proetale{{pro\text{-}\acute{e}tale}}
\def\etale{{\acute{e}tale}}
\def\QCoh{\mathit{QCoh}}
\def\Ker{\mathop{\mathrm{Ker}}}
\def\Im{\mathop{\mathrm{Im}}}
\def\Coker{\mathop{\mathrm{Coker}}}
\def\Coim{\mathop{\mathrm{Coim}}}

% Boxtimes
%
\DeclareMathSymbol{\boxtimes}{\mathbin}{AMSa}{"02}

%
% Macros for moduli stacks/spaces
%
\def\QCohstack{\mathcal{QC}\!\mathit{oh}}
\def\Cohstack{\mathcal{C}\!\mathit{oh}}
\def\Spacesstack{\mathcal{S}\!\mathit{paces}}
\def\Quotfunctor{\mathrm{Quot}}
\def\Hilbfunctor{\mathrm{Hilb}}
\def\Curvesstack{\mathcal{C}\!\mathit{urves}}
\def\Polarizedstack{\mathcal{P}\!\mathit{olarized}}
\def\Complexesstack{\mathcal{C}\!\mathit{omplexes}}
% \Pic is the operator that assigns to X its picard group, usage \Pic(X)
% \Picardstack_{X/B} denotes the Picard stack of X over B
% \Picardfunctor_{X/B} denotes the Picard functor of X over B
\def\Pic{\mathop{\mathrm{Pic}}\nolimits}
\def\Picardstack{\mathcal{P}\!\mathit{ic}}
\def\Picardfunctor{\mathrm{Pic}}
\def\Deformationcategory{\mathcal{D}\!\mathit{ef}}
\usepackage{cleveref}
\crefname{originallemma}{original lemma}{original lemmas}
\Crefname{originallemma}{Original lemma}{Original lemmas}
\makeatletter
\newlabel{lemma-generic-flatness-stratification-scheme}{{0H3Z}{source}{}{sourcepointer-1}{}}
\newlabel{lemma-generic-flatness-stratification-scheme@cref}{{[originallemma][0][]0H3Z}{[1][0][]source}{}{}{}}
\makeatother
\crefname{originalreference}{original reference}{original references}
\Crefname{originalreference}{Original reference}{Original references}
\makeatletter
\newlabel{more-algebra-lemma-check-tor-dimension-modulo-nilpotent}{{0H75}{source}{}{sourcepointer-2}{}}
\newlabel{more-algebra-lemma-check-tor-dimension-modulo-nilpotent@cref}{{[originalreference][0][]0H75}{[1][0][]source}{}{}{}}
\makeatother
\crefname{originalreference}{original reference}{original references}
\Crefname{originalreference}{Original reference}{Original references}
\makeatletter
\newlabel{more-algebra-lemma-f-flat}{{0H85}{source}{}{sourcepointer-3}{}}
\newlabel{more-algebra-lemma-f-flat@cref}{{[originalreference][0][]0H85}{[1][0][]source}{}{}{}}
\makeatother
\crefname{originalreference}{original reference}{original references}
\Crefname{originalreference}{Original reference}{Original references}
\makeatletter
\newlabel{perfect-lemma-affine-compare-bounded}{{06Z0}{source}{}{sourcepointer-4}{}}
\newlabel{perfect-lemma-affine-compare-bounded@cref}{{[originalreference][0][]06Z0}{[1][0][]source}{}{}{}}
\makeatother
\title{Testing flatness of a module in degree zero by pullback along a surjective finitely presented morphism}
\author{The Stacks Project Authors}
\date{stacks-project:more-morphisms.tex:local-snapshot-20261011; GNU FDL 1.2 or later}
\begin{document}
\maketitle
\noindent\textit{Editorial note.} Counted span is the author's own proof environment, source lines 15645-15698, immediately following the statement at 15632-15643.\par\medskip
\noindent\textit{Editorial note.} The statement carries the author's own attribution bracket 'Longke Tang'; it is retained verbatim inside the counted statement span.\par\medskip
\noindent\textit{Editorial note.} Three cross-chapter pointers (More on Algebra, Derived Categories of Schemes) plus one same-chapter pointer are recorded in assembly.external\_references.\par\medskip
\noindent\textit{Allowed standard prerequisites (editorial).}\begin{itemize}
\item Flatness versus tor amplitude in degree zero. More on Algebra, lemma more-algebra-lemma-f-flat, cited at proof line 15697. Converts the conclusion of the reduction (tor amplitude in [0,0]) into flatness of the module.
\item Checking tor dimension after reduction modulo a nilpotent. More on Algebra, lemma more-algebra-lemma-check-tor-dimension-modulo-nilpotent, cited in the proof. Allows the tor amplitude condition to be verified on the graded pieces of the stratification produced in the proof.
\item Affine comparison for bounded complexes. Derived Categories of Schemes, lemma perfect-lemma-affine-compare-bounded, cited in the proof. Carries the computation between the affine cover and the scheme level.
\item Stratification by flatness locus. chapter lemma lemma-generic-flatness-stratification-scheme, cited in the proof. Produces the stratification S' x S\_0 x ... x S\_t by the number of defining equations used in the induction.
\end{itemize}\medskip
\ifdefined\hypertarget\hypertarget{sourcepointer-1}{}\fi
\noindent\textit{Source reference.} Original lemma 0H3Z: sources/stacks-more-morphisms/pinned-original-source.tex, source lines 15615--15615. Navigation pointer only: the counted proof cites this statement by \textbackslash{}ref as a standing input; the deduction the counted proof has to make is carried out inside the counted span, and no part of that deduction is delegated to this pointer.\par\medskip
\ifdefined\hypertarget\hypertarget{sourcepointer-2}{}\fi
\noindent\textit{Source reference.} Original reference 0H75: focused-seeds/geometry/S9-morph-flatness-criterion-tang/ref-tags.txt, source lines 4--4. Navigation pointer only: the counted proof cites this statement by \textbackslash{}ref as a standing input; the deduction the counted proof has to make is carried out inside the counted span, and no part of that deduction is delegated to this pointer.\par\medskip
\ifdefined\hypertarget\hypertarget{sourcepointer-3}{}\fi
\noindent\textit{Source reference.} Original reference 0H85: focused-seeds/geometry/S9-morph-flatness-criterion-tang/ref-tags.txt, source lines 5--5. Navigation pointer only: the counted proof cites this statement by \textbackslash{}ref as a standing input; the deduction the counted proof has to make is carried out inside the counted span, and no part of that deduction is delegated to this pointer.\par\medskip
\ifdefined\hypertarget\hypertarget{sourcepointer-4}{}\fi
\noindent\textit{Source reference.} Original reference 06Z0: focused-seeds/geometry/S9-morph-flatness-criterion-tang/ref-tags.txt, source lines 6--6. Navigation pointer only: the counted proof cites this statement by \textbackslash{}ref as a standing input; the deduction the counted proof has to make is carried out inside the counted span, and no part of that deduction is delegated to this pointer.\par\medskip
\begin{lemma}[Longke Tang]
\label{lemma-flatness-criterion}
Let $f : X \to S$ be a surjective morphism of finite presentation.
Let $M$ be an object of $D_\QCoh(\mathcal{O}_S)$ such that
$H^i(M) = 0$ for $i > 0$. The following are equivalent
\begin{enumerate}
\item $M$ is isomorphic to a flat $\mathcal{O}_S$-module
placed in degree $0$,
\item $Lf^*M$ is isomorphic to a flat $\mathcal{O}_X$-module
placed in degree $0$.
\end{enumerate}
\end{lemma}

\begin{proof}
We omit the proof that (1) implies (2). Assume (2). To prove (1)
is a local question, hence we may assume that $S$ is affine.
Choose a stratification
$S \supset S_0 \supset S_1 \supset \ldots \supset S_t = \emptyset$
as in Lemma \ref{lemma-generic-flatness-stratification-scheme}.
Set $T_i = S_i \setminus S_{i - 1}$ and $Y_i = X \times_S T_i$.
Denote $g_i : T_i \to S$ and $h_i : Y_i \to X$ the inclusion morphisms.
Since $f$ is surjective, the morphisms $f_i : X_i \to T_i$
are surjective and flat. Since
$Lh_i^* \circ Lf^* = Lf_i^* \circ Lg_i^*$ we conclude from (2) that
$Lf_i^* Lg_i^*M$ is isomorphic to a flat $\mathcal{O}_{Y_i}$-module
placed in degree $0$. Since $f_i$ is flat and surjective, we see that
$Lg_i^*M$ is isomorphic to a flat $\mathcal{O}_{T_i}$-module
placed in degree $0$.

\medskip\noindent
We will prove this implies the result by induction on $t$.
We will use the following notation: $S = \Spec(A)$ and $M$
is the object of $D_\QCoh(\mathcal{O}_S)$ corresponding to
$N$ in $D(A)$, see Derived Categories of Schemes, Lemma
\ref{perfect-lemma-affine-compare-bounded}.

\medskip\noindent
Base case: $t = 1$. Write $S_0 = \Spec(A/I)$ where
$I$ is a finite generated ideal. Then $I$ is nilpotent as $S_0 = S$
set theoretically. The assumption is that $N \otimes_A^\mathbf{L} A/I$
has tor amplitude in $[0, 0]$. By More on Algebra, Lemma
\ref{more-algebra-lemma-check-tor-dimension-modulo-nilpotent}
the same is true for $N$.

\medskip\noindent
Induction step. Assume $t > 1$. Write $S_{t - 1} = \Spec(A/I)$
where $I = (f_1, \ldots, f_r)$ is a finitely generated ideal.
We will argue by induction on $r$. Observe that
$N_{f_r} = N \otimes_A^\mathbf{L} A_{f_r}$ in $D(A_{f_r})$
has tor amplitude in $[0, 0]$ by induction hypothesis
(because the stratification on the principal open $D(f_r)$ has length
at most $t - 1$). On the other hand, consider
$N' = N \otimes_A^\mathbf{L} A/f_rA$ in $D(A/f_rA)$.
Setting $S' = \Spec(A/f_rA)$ we obtain a stratification
$$
S' \supset S' \cap S_0 \supset S' \cap S_1 \supset \ldots \supset
S' \cap S_{t - 1} \supset
S' \cap S_t = \emptyset
$$
where $S' \cap S_{t - 1}$ is cut out by $r - 1$ equations in $S'$.
Similarly to before, the derived pullbacks of $M$ to the parts
$S' \cap S_i \setminus S' \cap S_{i + 1}$ are of tor amplitude
in $[0, 0]$. Whence by induction on $r$ we conclude that
$N'$ has tor amplitude in $[0, 0]$.
Then we finally conclude that $N$ has tor amplitude in $[0, 0]$
by More on Algebra, Lemma \ref{more-algebra-lemma-f-flat}.
\end{proof}

\end{document}
